\chapter{Chapter06 — Affine}
%%% generated by the blueprint pipeline from FABL.Chapter06.F₂Polynomials.Affine
\section{Overview}

Declarations of the Lean module \texttt{FABL.Chapter06.F₂Polynomials.Affine}.

\section{Declarations}

\begin{definition}[Affine function on the Boolean cube]
\lean{FABL.affineFunction}
\statementsource{Analysis-of-Boolean-Functions-by-Ryan-ODonnell}{pdf-ffbe9c7b3306-p038-b010}
\label{FABL.affineFunction}
\leanok
\uses{FABL.F₂BooleanFunction, FABL.F₂Cube, FABL.f₂DotProduct, FABL.𝔽₂}
Let $n$ be a natural number, let $bin_2$, and let $ain_2^n$. The affine function
with constant term $b$ and linear coefficients $a$ sends each $xin_2^n$ to $b$ plus
the dot product of $a$ and $x$.
\end{definition}

\begin{definition}[Coefficients of an affine function]
\lean{FABL.affineCoefficients}
\statementsource{Analysis-of-Boolean-Functions-by-Ryan-ODonnell}{pdf-ffbe9c7b3306-p038-b010}
\label{FABL.affineCoefficients}
\leanok
\uses{FABL.ANFCoefficients, FABL.F₂Cube, FABL.𝔽₂}
Let $n$ be a natural number, let $bin_2$, and let $ain_2^n$. Define coefficients on
subsets of $_n$ by assigning $b$ to the empty set, the sum of the coordinates $a_i$ for
$i$ in $S$ when $S$ has cardinality one, and $0$ otherwise.
\end{definition}

\begin{theorem}[Affine functions from their coefficients]
\lean{FABL.anfEval_affineCoefficients}
\label{FABL.anfEval_affineCoefficients}
\leanok
\uses{FABL.F₂Cube, FABL.affineCoefficients, FABL.affineFunction, FABL.anfEval, FABL.anfMonomial, FABL.f₂DotProduct, FABL.𝔽₂}
\difficulty{5}
Let $b\in\mathbb{F}_2$ and $a\in\mathbb{F}_2^n$. The algebraic normal form with constant
coefficient $b$, singleton coefficient at $\{i\}$ equal to $a_i$, and all other
coefficients zero evaluates at every $x\in\mathbb{F}_2^n$ to $b+\sum_{i=1}^n a_i x_i$.
\end{theorem}

\begin{theorem}[ANF coefficients of an affine function]
\lean{FABL.anfCoeff_affineFunction}
\label{FABL.anfCoeff_affineFunction}
\leanok
\uses{FABL.F₂Cube, FABL.affineCoefficients, FABL.affineFunction, FABL.anfCoeff, FABL.anfEval_affineCoefficients, FABL.anfEval_anfCoeff, FABL.anfEval_injective, FABL.𝔽₂}
\difficulty{3}
For $b\in\mathbb F_2$ and $a\in\mathbb F_2^n$, the coefficient function of $x\mapsto
b+a\cdot x$ has value $b$ on the empty set, value $\sum_{i\in S}a_i$ on each singleton
set $S$, and value $0$ on every other set of coordinates.
\end{theorem}

\begin{theorem}[Degree of an affine Boolean function]
\lean{FABL.functionAlgebraicDegree_affineFunction_le_one}
\label{FABL.functionAlgebraicDegree_affineFunction_le_one}
\leanok
\uses{FABL.F₂Cube, FABL.affineCoefficients, FABL.affineFunction, FABL.algebraicDegree_le_iff, FABL.anfCoeff_affineFunction, FABL.functionAlgebraicDegree, FABL.𝔽₂}
\difficulty{4}
For $b\in\mathbb{F}_2$ and $a$ in the $n$-dimensional cube over $\mathbb{F}_2$, the
affine Boolean function $x\mapsto b+a\cdot x$ has algebraic degree at most $1$, where
the dot product and addition are over $\mathbb{F}_2$.
\end{theorem}

\begin{theorem}[Degree-one functions are affine]
\lean{FABL.exists_affineFunction_of_functionAlgebraicDegree_le_one}
\label{FABL.exists_affineFunction_of_functionAlgebraicDegree_le_one}
\leanok
\uses{FABL.F₂BooleanFunction, FABL.F₂Cube, FABL.affineCoefficients, FABL.affineFunction, FABL.algebraicDegree_le_iff, FABL.anfCoeff, FABL.anfEval_affineCoefficients, FABL.anfEval_anfCoeff, FABL.functionAlgebraicDegree}
\difficulty{5}
Let $n$ be a nonnegative integer and let $f:\mathbb F_2^n\to\mathbb F_2$ have algebraic
degree at most one, with degree defined from its algebraic normal form coefficients.
Then there are $b\in\mathbb F_2$ and $a\in\mathbb F_2^n$ such that $f(x)=b+a\cdot x$ for
every $x\in\mathbb F_2^n$.
\end{theorem}

\begin{theorem}[Degree of an affine map coordinate]
\lean{FABL.functionAlgebraicDegree_affineMap_coordinate_le_one}
\label{FABL.functionAlgebraicDegree_affineMap_coordinate_le_one}
\leanok
\uses{FABL.F₂Cube, FABL.IsF₂Linear, FABL.affineFunction, FABL.functionAlgebraicDegree, FABL.functionAlgebraicDegree_affineFunction_le_one, FABL.f₂DotProduct, FABL.isF₂Linear_iff_exists_dotProduct, FABL.𝔽₂}
\difficulty{4}
Let $L$ be an affine map from the $n$-dimensional cube over $\mathbb{F}_2$ to itself,
and let $i$ be a coordinate index. The Boolean function that sends $x$ to the $i$th
coordinate of $L(x)$ has algebraic degree at most $1$.
\end{theorem}

\begin{theorem}[Degree of an affine substitution in a monomial]
\lean{FABL.functionAlgebraicDegree_anfMonomial_comp_affineMap_le_card}
\label{FABL.functionAlgebraicDegree_anfMonomial_comp_affineMap_le_card}
\leanok
\uses{FABL.F₂Cube, FABL.anfMonomial, FABL.functionAlgebraicDegree, FABL.functionAlgebraicDegree_affineMap_coordinate_le_one, FABL.functionAlgebraicDegree_finset_prod_le, FABL.𝔽₂}
\difficulty{3}
Let $L$ be an affine map from the $n$-dimensional cube over $\mathbb{F}_2$ to itself,
and let $S$ be a subset of the $n$ coordinate indices. The function $x\mapsto\prod_{i\in
S}(L(x))_i$ has algebraic degree at most $|S|$, with the product taken in
$\mathbb{F}_2$.
\end{theorem}

\begin{theorem}[Degree under affine composition]
\lean{FABL.functionAlgebraicDegree_comp_affineMap_le}
\label{FABL.functionAlgebraicDegree_comp_affineMap_le}
\leanok
\uses{FABL.F₂BooleanFunction, FABL.F₂Cube, FABL.algebraicDegree_le_iff, FABL.anfCoeff, FABL.anfEval_anfCoeff, FABL.anfMonomial, FABL.functionAlgebraicDegree, FABL.functionAlgebraicDegree_anfMonomial_comp_affineMap_le_card, FABL.functionAlgebraicDegree_finset_sum_le, FABL.functionAlgebraicDegree_zero, FABL.𝔽₂}
\difficulty{4}
Let $n$ be a natural number, let $f$ be a function from the $n$-dimensional Boolean cube
over the two-element field to that field, and let $L$ be any affine map from this cube
to itself. The algebraic degree of $f$ composed with $L$ is at most the algebraic degree
of $f$.
\end{theorem}

\begin{theorem}[Affine invariance of algebraic degree]
\lean{FABL.functionAlgebraicDegree_comp_affineEquiv}
\label{FABL.functionAlgebraicDegree_comp_affineEquiv}
\leanok
\uses{FABL.F₂BooleanFunction, FABL.F₂Cube, FABL.functionAlgebraicDegree, FABL.functionAlgebraicDegree_comp_affineMap_le, FABL.𝔽₂}
\difficulty{3}
Let $n$ be a natural number, let $f$ be a function from the $n$-dimensional Boolean cube
over the two-element field to that field, and let $L$ be an affine bijection of this
cube. The algebraic degree of $f$ composed with $L$ equals the algebraic degree of $f$.
\end{theorem}
